\documentclass[addpoints,10pt]{exam}
\usepackage[utf8]{inputenc}
\usepackage[T1]{fontenc}
\usepackage[spanish]{babel}
\usepackage{times}
\usepackage{enumerate}
\usepackage{graphicx}
\usepackage{stfloats}
\usepackage{multirow}
\usepackage{array} 
\usepackage{vhistory}
\usepackage{tcolorbox}
\usepackage{fancyvrb}
\usepackage[lighttt]{lmodern}
\usepackage{listings}
\lstset{basicstyle=\small\ttfamily,
	stringstyle=\ttfamily,
	commentstyle=\color{gray45},
	showstringspaces = false,
	keywordstyle=\bfseries}
\usepackage{multicol,caption,capt-of}
%multi-row
\usepackage{multirow}
\usepackage[left=1.5cm,right=1.5cm,top=1.5cm,bottom=1.5cm]{geometry}

\begin{document}
	
	%Encabezados y pie de página
	%----------------------------------------------------------------------------------------
	\pagestyle{headandfoot}
	\runningheadrule
	\firstpageheader{Sistemas Computacionales}{Evaluación Unidad 7 - GPU}{Mayo 11 de 2026}
	\runningheader{Sistemas Computacionales}
	{Evaluación Unidad 7, Page \thepage\ of \numpages}
	{Mayo 11 de 2026}
	\runningfooter{Escuela de Ingenierías}{UPB}{Ingeniería en Diseño de Entretenimiento Digital}
	
	\begin{center}
	\begin{tabular}{|m{4cm}|m{8cm}|m{5cm}|}
		\hline 
		\multirow{5}{*}{\includegraphics[width=1\linewidth]{Logo-UPB-2022.pdf}}
		&\multicolumn{2}{c|}{\textbf{Universidad Pontificia Bolivariana - Sede  Medellín}}		\\ \cline{2-3}
		&Curso: \textbf{Sistemas Computacionales}		&Duración: 90 minutos				\\ \cline{2-3}
		&Preparada por:\textbf{ Henry Andrade, IEo, Ph.D.}	&Evaluador:							\\ \cline{2-3}
		&Ing. en Diseño de Entretenimiento Digital	&Prueba: Convencional				\\ \cline{2-3}
		&\multicolumn{2}{l|}{Estudiante:				}										\\ 
		 \\ \hline
	\end{tabular} 
\end{center}

	%----------------------------------------------------------------------------------------
	%Cuadro con instrucciones
	%----------------------------------------------------------------------------------------		
	\tcbset{colback=red!5!white, colframe=red!75!black, fonttitle=\bfseries}
	\begin{tcolorbox}[title=Nota Importante]
		Lea cuidadosamente cada una de las preguntas antes de contestar. Durante el examen, no se autoriza el acceso a Internet por medio de ningún dispositivo. Cualquier intento de copia o fraude dará lugar a la anulación del examen.
	\end{tcolorbox}
	%----------------------------------------------------------------------------------------
	
	\qformat{\textbf{Pregunta \thequestion} \dotfill (\thepoints\ puntos)}
	\boxedpoints
	
	\begin{questions}
		
	% ---------------------------------
	% PREGUNTAS DE SELECCIÓN MÚLTIPLE (20 pts)
	% ---------------------------------
	\part{Preguntas de Selección Múltiple (20 puntos)}
	
	\question[4]
	¿Cuál es la función principal de la biblioteca \textbf{GLAD} en un proyecto de OpenGL moderno?
	\begin{oneparchoices}
		\choice Crear y gestionar ventanas y contextos de OpenGL.
		\choice Cargar dinámicamente las funciones de OpenGL desde los drivers de la GPU en tiempo de ejecución.
		\choice Proveer funciones matemáticas para transformaciones de vectores y matrices.
		\choice Gestionar la compilación y el enlace de los shaders (vertex y fragment).
	\end{oneparchoices}
	\begin{solution}
	La respuesta correcta es la B. GLAD es un cargador de funciones que resuelve el problema de que las funciones de OpenGL (versiones > 1.1) no están directamente disponibles, sino que deben ser obtenidas desde el driver de la tarjeta gráfica. GLFW (A) se encarga de las ventanas, GLM (C) de las matemáticas y el propio código de la aplicación (D) gestiona los shaders.
	\end{solution}

	\question[4]
	En el pipeline de gráficos de OpenGL, ¿qué información fundamental procesa el \textbf{Vertex Shader}?
	\begin{oneparchoices}
		\choice El color final de cada píxel que se mostrará en pantalla.
		\choice Los datos de entrada por vértice (como la posición) y calcula su posición final en coordenadas de recorte (clip space).
		\choice La visibilidad de los objetos, descartando aquellos que están fuera del campo de visión.
		\choice La textura que se aplicará a cada triángulo de la malla.
	\end{oneparchoices}
	\begin{solution}
	La respuesta correcta es la B. El Vertex Shader se ejecuta una vez por cada vértice y su responsabilidad principal es tomar los atributos de ese vértice (especialmente la posición en `aPos`) y transformarlos a coordenadas de recorte, asignando el resultado a `gl_Position`. El Fragment Shader (A) calcula el color del píxel, el frustum culling (C) es un paso previo, y el mapeo de texturas (D) ocurre mayormente en el Fragment Shader.
	\end{solution}

	\question[4]
	¿Qué representa un \textbf{Vertex Array Object (VAO)} en OpenGL?
	\begin{oneparchoices}
		\choice Un bloque de memoria en la GPU que almacena únicamente las coordenadas de los vértices.
		\choice Un objeto que encapsula toda la configuración del estado de los atributos de vértice y los VBOs asociados.
		\choice Un programa que se ejecuta en la GPU para determinar el color de cada fragmento.
		\choice Una variable global en un shader que permite pasar datos desde la CPU a la GPU.
	\end{oneparchoices}
	\begin{solution}
	La respuesta correcta es la B. El VAO actúa como un contenedor que "recuerda" todas las configuraciones realizadas con `glVertexAttribPointer`, `glEnableVertexAttribArray` y el VBO que estaba enlazado en ese momento. Esto simplifica enormemente el proceso de dibujo, ya que solo se necesita enlazar el VAO para restaurar todo el estado de los vértices. Un VBO (A) almacena los datos, un Fragment Shader (C) colorea fragmentos y una variable `uniform` (D) pasa datos globales.
	\end{solution}

	\question[4]
	¿Cuál es el propósito de una variable \texttt{uniform} en GLSL?
	\begin{oneparchoices}
		\choice Proporcionar datos diferentes para cada vértice, como la posición o el color.
		\choice Pasar datos desde la CPU a la GPU que permanecen constantes para todos los vértices y fragmentos en un mismo `draw call`.
		\choice Almacenar el color de salida final de un Fragment Shader.
		\choice Interpolar valores (como el color) a través de la superficie de un triángulo durante la rasterización.
	\end{oneparchoices}
	\begin{solution}
	La respuesta correcta es la B. A diferencia de los atributos de vértice, los `uniforms` son globales para una ejecución de shader. Son ideales para pasar datos como una matriz de transformación, el tiempo transcurrido o un color base que afectará a todo el objeto de la misma manera. Los atributos (`in`) proporcionan datos por vértice (A), las variables de salida (`out`) como `FragColor` almacenan el resultado (C), y la interpolación (D) es un proceso automático del pipeline.
	\end{solution}

	\question[4]
	¿Qué hace la función \texttt{glfwSwapBuffers(window)} en el ciclo de renderizado?
	\begin{oneparchoices}
		\choice Limpia el framebuffer con un color de fondo antes de dibujar un nuevo cuadro.
		\choice Intercambia el buffer frontal (que se está mostrando en pantalla) con el buffer trasero (donde se acaba de dibujar la escena).
		\choice Copia los datos de los vértices desde la RAM de la CPU a la VRAM de la GPU.
		\choice Activa la sincronización vertical (VSync) para evitar el `tearing`.
	\end{oneparchoices}
	\begin{solution}
	La respuesta correcta es la B. Esta función es clave en la técnica de doble búfer. Mientras la pantalla muestra el `front buffer`, la GPU dibuja la siguiente escena en un `back buffer` oculto. `glfwSwapBuffers` presenta el cuadro recién completado intercambiando los roles de estos búferes, lo que previene que se vean artefactos de dibujo a medio terminar. `glClear` (A) limpia el buffer, `glBufferData` (C) copia datos, y `glfwSwapInterval` (D) configura VSync.
	\end{solution}
	
	% --------------------------
	% PREGUNTAS ABIERTAS (30 pts)
	% --------------------------
	\part{Preguntas Abiertas (30 puntos)}
	
	\question[10]
	Explica la relación y el propósito de los siguientes tres componentes en OpenGL: \textbf{VBO (Vertex Buffer Object)}, \textbf{VAO (Vertex Array Object)} y la función \texttt{glVertexAttribPointer()}. Describe cómo trabajan juntos para enviar datos de vértices a un Vertex Shader.
	
	\begin{solution}[\stretch{2in}]
	\textbf{Rúbrica de Evaluación (10 puntos):}
	\begin{itemize}
		\item \textbf{(3 pts)} Define correctamente el VBO como un búfer de memoria en la GPU que almacena los datos brutos de los vértices (posiciones, colores, etc.).
		\item \textbf{(3 pts)} Define correctamente el VAO como un objeto de estado que almacena la configuración de cómo OpenGL debe interpretar los datos de los VBOs. Menciona que "recuerda" los enlaces de `glVertexAttribPointer`.
		\item \textbf{(4 pts)} Explica claramente la función de `glVertexAttribPointer` como el "manual de instrucciones" que le dice a OpenGL cómo leer el VBO actualmente enlazado: dónde empieza cada atributo dentro del bloque de datos de un vértice (offset), cuánto espacio ocupa (size), y cuánto hay que saltar para encontrar el mismo atributo en el siguiente vértice (stride). Vincula el primer parámetro (`index`) con el `layout(location = ...)` en el shader.
	\end{itemize}
	\end{solution}

	\question[10]
	Describe el concepto de \textbf{pipeline de gráficos programable} introducido en el OpenGL moderno. Compara las responsabilidades del \textbf{Vertex Shader} y el \textbf{Fragment Shader} dentro de este pipeline y explica qué tipo de efecto visual podrías implementar en cada uno.
	
	\begin{solution}[\stretch{2.5in}]
	\textbf{Rúbrica de Evaluación (10 puntos):}
	\begin{itemize}
		\item \textbf{(3 pts)} Explica que el pipeline programable significa que ciertas etapas clave del renderizado (antes de función fija) ahora son controladas por el programador mediante pequeños programas llamados shaders, que se ejecutan directamente en la GPU.
		\item \textbf{(3 pts)} Describe la responsabilidad del Vertex Shader: opera por vértice, principalmente para transformar sus coordenadas espaciales. Menciona un efecto adecuado como mover o escalar el objeto, o aplicar una deformación (ej. efecto de onda).
		\item \textbf{(3 pts)} Describe la responsabilidad del Fragment Shader: opera por "fragmento" (potencial píxel), principalmente para calcular su color final. Menciona un efecto adecuado como aplicar un color sólido, un degradado, leer de una textura, o calcular iluminación básica.
		\item \textbf{(1 pt)} La explicación general es coherente y demuestra una clara diferenciación entre las dos etapas.
	\end{itemize}
	\end{solution}

	\question[10]
	En la Actividad 6, se propone hacer que el color de un triángulo cambie con el tiempo usando la función \texttt{glfwGetTime()} y una variable \texttt{uniform}. Describe el proceso completo para implementar este efecto. Tu respuesta debe incluir:
	\begin{enumerate}
		\item La modificación necesaria en el Fragment Shader.
		\item Los pasos a seguir en el código C++ (obtener la ubicación del uniform y actualizar su valor en el ciclo de renderizado).
	\end{enumerate}
	
	\begin{solution}[\stretch{3in}]
	\textbf{Rúbrica de Evaluación (10 puntos):}
	\begin{itemize}
		\item \textbf{(4 pts)} \textbf{Modificación del Fragment Shader:} Escribe o describe correctamente un Fragment Shader que declara una variable `uniform float u_time;`. Luego, en `main()`, utiliza `u_time` con una función periódica (como `sin()` o `cos()`) para calcular un valor de color dinámico. Por ejemplo: `FragColor = vec4((sin(u_time) + 1.0) / 2.0, 0.5, 0.2, 1.0);`.
		\item \textbf{(2 pts)} \textbf{Obtener Ubicación en C++:} Explica que, después de crear y enlazar el programa de shaders, se debe llamar a `glGetUniformLocation(shaderProg, "u_time");` para obtener un ID entero que representa la ubicación de esa variable `uniform`. Este paso se hace una sola vez, antes del ciclo principal.
		\item \textbf{(4 pts)} \textbf{Actualizar en el Ciclo de Renderizado:} Explica que, dentro del bucle `while`, en cada fotograma, se debe obtener el tiempo actual con `glfwGetTime()`. Luego, justo antes de la llamada a `glDrawArrays`, se debe actualizar el valor del `uniform` usando la ubicación obtenida y el tiempo actual, con una llamada como `glUniform1f(timeLocation, (float)glfwGetTime());`.
	\end{itemize}
	\end{solution}

	\end{questions}
	
	\begin{center}
		\gradetable[h][questions]
	\end{center}
	
\end{document}